\chapter{Safety Requirements}

interface (fixed hole/rail pattern) documented for reuse. FR-33       The mounting area shall provide a power and data                 S           I,T take-off point so future modules can draw power and exchange data with the base. FR-34       A representative dummy module shall be mountable                 S            D and removable using hand tools within 10 minutes.

5 Safety Requirements These are the minimum safety requirements for V1. As a prototype, V1 operates in a controlled, restricted test area under supervision (see Constraints, Section 6); it is not yet expected to hold full third-party safety certification. The design should nonetheless follow the intent of the rel- evant standards ISO 3691-4 (driverless industrial trucks), ANSI/RIA R15.08 (industrial mobile robots), ISO 13849 (safety-related control), and IEC 60204 so later versions can be certified without redesign.

ID          Requirement description                                        Pri.         Verif. SR-01       The robot shall provide at least one clearly marked,            M             T easily reachable hardware emergency-stop (E-stop) that removes drive power. SR-02       Activating the E-stop shall bring the robot to a                M             T complete stop within 0.5 s in all operating modes. SR-03       On loss of power or critical fault, the braking system          M             T shall fail safe (robot stops and holds position).

ID          Requirement description                                         Pri.          Verif. SR-04       The robot shall automatically slow and stop for                     M           T obstacles detected in its path before contact (links to FR-15). SR-05       The robot shall limit speed near detected                           M           T obstacles/people to $\leq$0.3 m/s and shall never exceed the 1.5 m/s hard limit. SR-06       The robot shall provide audible and visual indication               M           D when powered/moving (e.g. beacon and buzzer). SR-07       The robot shall remain stable (no tip-over) within its              M          A,T rated payload, speed, and maneuver envelope. SR-08       Battery and power electronics shall be protected against            M          I,T over-current, over-temperature, and short circuit (links to FR-23/24). SR-09       Exposed moving parts and pinch points shall be                      S            I guarded or eliminated within the V1 envelope. SR-10       The robot shall provide contact-based protection (e.g.              S           T bumper or contact sensor) as a last line of defence

